\documentclass{cap2026}
% Use \documentclass[french]{cap2026} if you want to load the french version of the template.
% Note that if you choose to use the french option, you might encounter some unexpected behaviors, for example if you are using colons (:) in your labels.
% Remove the french option if you want to load the english version of the template.

% The following packages are already loaded:
% authblk
% hyperref
% natbib - Use the option no-natbib to avoid loading it (\documentclass[no-natbib]{cap2026})
% geometry with the options width=17cm and height=24cm
% amsthm
% and some environments are already defined and automatically switch between french and english: theorem, definition, lemma, corollary, proposition. 

\usepackage{fancyvrb}
\usepackage{amsfonts} 
\usepackage{amssymb}
\usepackage{amsmath}
\usepackage{booktabs} % for better tables
\usepackage{cleveref}

\title{Example of an article in the CAp'2026 format}
\author[1]{Author A\thanks{If you really want to include your email
    or web page...: A.A@univ-x.fr}}
\author[2]{Author B}
\affil[1]{University X, CNRS}
\affil[2]{University Y, CNRS and Inria}


\begin{document}
\maketitle

\begin{abstract}
  This style is very simple. You can add keywords below the abstract using the \texttt{keywords} command. The number of pages is limited to 8 or 12 pages (depending on the type of paper), excluding references.
  
\medskip

\keywords{Style, \LaTeX, CAp}
\end{abstract}


\section{The header}
\label{seclentete}

Nothing special to mention; this year we are opting for a single-column style. The \texttt{french} option automatically switches the language of the paper to French. Some packages are already loaded by default by the style:

\begin{Verbatim}
\usepackage{authblk}
\usepackage{hyperref}
\usepackage{natbib}
\usepackage[width=17cm,height=24cm]{geometry}
\usepackage{amsthm}
\end{Verbatim}

The \texttt{authblk} package is used for affiliations. The \texttt{hyperref} package makes navigation in the document easier. The \texttt{natbib} package is used by default to manage references but can, if needed, be disabled using the \texttt{no-natbib} option of the class. The \texttt{geometry} package manages the document margins. They must not be modified. Finally, the \texttt{amsthm} package manages theorem environments (the \texttt{theorem}, \texttt{definition}, \texttt{lemma}, \texttt{proposition}, and \texttt{corollary} environments are already available).

You can add any packages or macros you want, for example:
\begin{Verbatim}
  \usepackage{fancyvrb}
  \usepackage{amsfonts} 
  \usepackage{amssymb}
  \usepackage{amsmath}
\end{Verbatim}

\section{The body}
\label{secle-corps}

The \texttt{theorem}, \texttt{definition}, \texttt{lemma}, \texttt{proposition}, and \texttt{corollary} environments are already available and automatically switch language with the \texttt{french} option of the class.

\begin{theorem}[Bla]
  This is a nice theorem.
\end{theorem}

Bla bla bla bla bla bla bla bla bla bla bla bla bla bla bla bla bla
bla bla bla bla bla bla bla bla bla bla bla bla bla bla bla bla bla
bla bla bla bla bla bla bla bla bla bla bla bla bla bla bla bla bla
bla bla bla bla bla bla bla bla bla bla bla bla bla bla bla bla bla
bla bla bla bla bla bla bla bla bla bla bla bla bla bla bla bla bla
bla bla bla bla bla bla bla bla bla bla bla bla bla bla bla bla bla
bla bla bla bla bla bla bla bla bla bla


\section{Figures}
\label{secles-figures}
You can add tables (see \cref{tabtab}) and figures (see \cref{figcode}).

\begin{table}
  \centering  
  \begin{tabular}{llll}
    \toprule
    \textbf{Column 1} & \textbf{Column 2} & \textbf{Column 3} &
    \textbf{Column 4} \\ \midrule
    1 & 2 & 3 & 4 \\ 
    1 & 2 & 3 & 4 \\ 
    1 & 2 & 3 & 4 \\ 
    1 & 2 & 3 & 4 \\ \bottomrule
  \end{tabular}
  \caption{An example table.\label{tabtab}}
\end{table}

\begin{SaveVerbatim}{codefigure}
\begin{table}
  \centering  
  \begin{tabular}{llll}
    \toprule
    \textbf{Column 1} & \textbf{Column 2} & \textbf{Column 3} &
    \textbf{Column 4} \\ \midrule
    1 & 2 & 3 & 4 \\ 
    1 & 2 & 3 & 4 \\ 
    1 & 2 & 3 & 4 \\ 
    1 & 2 & 3 & 4 \\ \bottomrule
  \end{tabular}
  \caption{An example table.\label{tabtab}}
\end{table}
\end{SaveVerbatim}

\begin{figure}
  \centering
  \UseVerbatim[frame=single]{codefigure}
  \caption{An example figure.\label{figcode}}
\end{figure}

Bla bla bla bla bla bla bla bla bla bla bla bla bla bla bla bla bla
bla bla bla bla bla bla bla bla bla bla bla bla bla bla bla bla bla
bla bla bla bla bla bla bla bla bla bla bla bla bla bla bla bla bla
bla bla bla bla bla bla bla bla bla bla bla bla bla bla bla bla bla
bla bla bla bla bla bla bla bla bla bla bla bla bla bla bla bla bla
bla bla bla bla bla bla bla bla bla bla bla bla bla bla bla bla bla
bla bla bla bla bla bla bla bla bla bla



\section{Bibliography}
\label{secla-bibliographie}

We recommend using \texttt{natbib} with the \texttt{apalike} style, but we do not require it. The advantage of \texttt{natbib} is that it automatically manages citations in the text, for example~\citet{paper1}, and in parentheses, for example~\citep{DD14a}.

Bla bla bla bla bla bla bla bla bla bla bla bla bla bla bla bla bla
bla bla bla bla bla bla bla bla bla bla bla bla bla bla bla bla bla
bla bla bla bla bla bla bla bla bla bla bla bla bla bla bla bla bla
bla bla bla bla bla bla bla bla bla bla bla bla bla bla bla bla bla
bla bla bla bla bla bla bla bla bla bla bla bla bla bla bla bla bla
bla bla bla bla bla bla bla bla bla bla bla bla bla bla bla bla bla
bla bla bla bla bla bla bla bla bla bla

\section{Without \LaTeX}
\label{secsans-latex}

If you do not use \LaTeX, you can follow the dimensions given at the beginning of this document and try to find a readable font.


\bibliographystyle{apalike}
\bibliography{cap2026}

\end{document}

%%% Local Variables: 
%%% mode: latex
%%% TeX-master: t
%%% End: 
